% META: {"id":"1SPE-EXPO-EX-006","chapitre":"1SPE-EXPONENTIELLE","type_objet":"exercice","capacites":["1SPE-EXPONENTIELLE-C1","1SPE-EXPONENTIELLE-C3"],"capacites_codes":["C1","C3"],"methodes":["M2"],"parcours":2,"competences":["calculer","representer"],"duree_min":8,"mode_creation":"ex_nihilo","sources_inspiration":[],"fichier_tex":"chapitres/1SPE-EXPONENTIELLE/exercices/1SPE-EXPO-EX-006.tex","corrige_tex":"chapitres/1SPE-EXPONENTIELLE/corriges/1SPE-EXPO-CO-006.tex","status":"generated"}
% BEGIN-VERIFY
% from sympy import *
% x = symbols('x')
% f = exp(x)
% # Tangent at x=0: y = f(0) + f'(0)*(x-0) = 1 + x
% f0 = f.subs(x, 0)
% fp0 = f.diff(x).subs(x, 0)
% tangent = f0 + fp0 * x
% assert tangent == 1 + x, "tangent at 0 is y = x + 1"
% assert f0 == 1, "f(0) = 1"
% assert fp0 == 1, "f'(0) = 1"
% END-VERIFY
\begin{exercice}{1SPE-EXPO-EX-006}{2}{8}
On considère la courbe représentative $\mathcal{C}$ de la fonction exponentielle dans un repère orthonormé.
\begin{enumerate}
\item Rappeler la valeur de $\exp(0)$ et $\exp'(0)$.
\item Déterminer l'équation de la tangente $T$ à la courbe $\mathcal{C}$ au point d'abscisse $0$.
\item La tangente passe-t-elle par l'origine du repère ? Justifier.
\item Étudier les variations de $h(x)=\mathrm{e}^x-(x+1)$, puis en déduire la position de la courbe $\mathcal{C}$ par rapport à la tangente $T$.
\end{enumerate}
\end{exercice}
